\section{関連研究} \label{sec2}
% \subsection{Retrieval-Augmented Generation (RAG)} \label{sec2.1}
% \subsection{適応的なRAG} \label{sec2.2}
質問の特徴に適応して検索方法やハイパーパラメータを選択する，適応的なRAGに関する研究が進展している~\cite{Asai23,fu24,Tang25,Rishi25,Zhang25}．
一般に，複雑な質問にはグラフベースRAG等の高精度な手法が有効である一方，単純な質問に複雑な処理を適用すると，不要な検索や計算によりコストが増大する．
一方，複雑な質問をNaiveRAGのような単純な手法や非検索で処理した場合には，応答品質が低下する~\cite{Jeong24}． 
適応的なRAGは，質問の複雑さに応じて検索方法やハイパーパラメータを選択し，応答品質の向上を目指す．

既存研究の多くは，LLMによる自己判定やオンライン強化学習に基づいている．
Self-RAG~\cite{Asai23}は，RAGの検索方法をLLMが自己評価しながら生成する手法である．
AutoRAG-HP~\cite{fu24}は，RAGの性能の要因となる複数のハイパーパラメータをオンライン強化学習で最適化する手法である．
Tangら~\cite{Tang25}は，グラフベースRAGにおいて，質問の特性に応じて最適な検索手法を自動的に選択する手法を提案した．
また，HyPA-RAG~\cite{Rishi25}は，質問の特性に応じて複数のハイパーパラメータを動的に最適化する手法である．
AdaGCRAG~\cite{Zhang25}は，小規模言語モデル環境においても高い性能と効率性を両立させる手法である．
これらの適応的なRAGは，検索方法の判断やバンディットによるオンライン最適化といった多様なアプローチにより実現されている．

これに対し，本研究では，質問と文書の多面的な特徴量および過去の評価結果に基づき，候補となるRAG構成をRMにより直接順位付けする．
これにより，精度とコストのトレードオフを同時に考慮したRAG構成の自動選択を実現する．
